\documentclass[UTF8,fontset=fandol,aspectratio=169,11pt]{ctexbeamer}
% Shared Beamer preamble for Big Data and Management Decision-Making slides.

\usetheme{Madrid}
\usecolortheme{default}
\usefonttheme{professionalfonts}

\usepackage{amsmath,amssymb,mathtools}
\usepackage{booktabs,array}
\usepackage{tikz}
\usepackage{graphicx}
\usepackage{listings}
\usepackage{xcolor}
\usetikzlibrary{arrows.meta,positioning,shapes.geometric}

\newcommand{\BDMFandolPath}{/usr/local/texlive/2025/texmf-dist/fonts/opentype/public/fandol/}
\setCJKmainfont[
  Path=\BDMFandolPath,
  UprightFont=FandolSong-Regular.otf,
  BoldFont=FandolSong-Regular.otf,
  ItalicFont=FandolKai-Regular.otf,
  BoldItalicFont=FandolKai-Regular.otf
]{FandolSong-Regular.otf}
\setCJKsansfont[
  Path=\BDMFandolPath,
  UprightFont=FandolSong-Regular.otf,
  BoldFont=FandolSong-Regular.otf
]{FandolSong-Regular.otf}
\setCJKmonofont[
  Path=\BDMFandolPath,
  UprightFont=FandolSong-Regular.otf,
  BoldFont=FandolSong-Regular.otf
]{FandolSong-Regular.otf}
\setmonofont{Menlo}

\definecolor{AcademicBlue}{RGB}{22,44,73}
\definecolor{AcademicTeal}{RGB}{22,119,119}
\definecolor{AcademicGray}{RGB}{76,84,95}
\definecolor{AcademicLight}{RGB}{245,247,250}
\definecolor{AcademicAmber}{RGB}{181,111,35}
\definecolor{CodeBg}{RGB}{248,248,248}

\setbeamercolor{structure}{fg=AcademicBlue}
\setbeamercolor{frametitle}{fg=white,bg=AcademicBlue}
\setbeamercolor{title}{fg=white,bg=AcademicBlue}
\setbeamercolor{block title}{fg=white,bg=AcademicBlue}
\setbeamercolor{block body}{bg=AcademicLight,fg=black}
\setbeamertemplate{navigation symbols}{}
\setbeamertemplate{footline}[frame number]
\setbeamertemplate{itemize item}{\raisebox{0.2ex}{\(\blacktriangleright\)}}
\setbeamerfont{title}{series=\mdseries}
\setbeamerfont{subtitle}{series=\mdseries}
\setbeamerfont{frametitle}{series=\mdseries}
\setbeamerfont{block title}{series=\mdseries}
\setbeamerfont{section in toc}{series=\mdseries}

\lstdefinestyle{pythonstyle}{
  basicstyle=\ttfamily\scriptsize,
  backgroundcolor=\color{CodeBg},
  keywordstyle=\color{AcademicBlue},
  commentstyle=\color{AcademicGray},
  stringstyle=\color{AcademicAmber},
  showstringspaces=false,
  breaklines=true,
  frame=single,
  rulecolor=\color{black!20},
  columns=fullflexible
}

\lstdefinestyle{sqlstyle}{
  language=SQL,
  basicstyle=\ttfamily\scriptsize,
  backgroundcolor=\color{CodeBg},
  keywordstyle=\color{AcademicBlue},
  commentstyle=\color{AcademicGray},
  stringstyle=\color{AcademicAmber},
  showstringspaces=false,
  breaklines=true,
  frame=single,
  rulecolor=\color{black!20},
  columns=fullflexible
}

\lstset{style=pythonstyle}

\hypersetup{
  colorlinks=true,
  linkcolor=AcademicBlue,
  urlcolor=AcademicTeal,
  pdfauthor={大数据与管理决策基础},
  pdfsubject={大数据与管理决策基础课程课件}
}


\title[流程数据与过程挖掘]{第10讲：流程数据与过程挖掘}
\subtitle{从事件日志到变体、瓶颈、返工与合规异常}
\author{《大数据与管理决策基础》}
\date{}

\begin{document}

\begin{frame}
  \titlepage
  \begin{center}
    \small 核心命题：流程不是静态图，而是由 case、activity 和 timestamp 构成的可计算行为证据。
  \end{center}
\end{frame}

\begin{frame}{本讲学习目标}
\begin{enumerate}
  \item 理解事件日志的基本数据模型；
  \item 解释 case、activity、timestamp、trace 和 variant；
  \item 构造 directly-follows graph；
  \item 计算吞吐时间、边耗时和活动频次；
  \item 识别返工、瓶颈和合规异常；
  \item 将过程挖掘结果转化为流程改进建议。
\end{enumerate}
\end{frame}

\begin{frame}{课程主线中的第10周}
\begin{center}
\resizebox{0.96\linewidth}{!}{%
\begin{tikzpicture}[
  node/.style={draw=AcademicBlue!65, fill=AcademicLight, rounded corners=1pt,
               minimum height=0.85cm, minimum width=2.05cm, align=center},
  active/.style={draw=AcademicBlue, fill=AcademicBlue, text=white,
               minimum height=0.9cm, minimum width=2.25cm, align=center},
  arrow/.style={-{Stealth[length=2mm]}, thick, AcademicTeal}
]
\node[node] (decision) {优化决策};
\node[node, right=0.65cm of decision] (execute) {流程执行};
\node[active, right=0.65cm of execute] (log) {事件日志};
\node[node, right=0.65cm of log] (improve) {流程改进};
\node[node, right=0.65cm of improve] (feedback) {反馈治理};
\draw[arrow] (decision) -- (execute);
\draw[arrow] (execute) -- (log);
\draw[arrow] (log) -- (improve);
\draw[arrow] (improve) -- (feedback);
\end{tikzpicture}
}
\end{center}
\small
优化方案只有进入流程执行后，才会暴露等待、返工、跳步和责任边界。
\end{frame}

\begin{frame}{事件日志}
\begin{center}
\begin{tabular}{p{0.24\linewidth}p{0.58\linewidth}}
\toprule
字段 & 含义\\
\midrule
case id & 一个流程实例，如订单、工单、申请。\\
activity & 实际发生的业务活动。\\
timestamp & 活动发生或完成时间。\\
resource & 执行人、系统、部门或设备。\\
attributes & 金额、区域、客户、产品等上下文。\\
\bottomrule
\end{tabular}
\end{center}
\[
  \sigma_c=(a_{c1},a_{c2},\ldots,a_{cm_c}).
\]
\small
按时间排序后的活动序列称为 trace。
\end{frame}

\begin{frame}{过程挖掘三类任务}
\begin{center}
\begin{tabular}{p{0.22\linewidth}p{0.34\linewidth}p{0.30\linewidth}}
\toprule
任务 & 目标 & 管理问题\\
\midrule
过程发现 & 从日志构造流程模型。 & 实际流程怎样运行？\\
合规检查 & 比较日志与参考流程。 & 是否跳步、反序或缺失？\\
流程增强 & 叠加时间、资源、成本。 & 哪里慢、贵、易返工？\\
\bottomrule
\end{tabular}
\end{center}
\end{frame}

\begin{frame}{Directly-Follows Graph}
若活动 \(a\) 在某个 case 中紧接活动 \(b\)，记为：
\[
  a\rightarrow b.
\]
\begin{itemize}
  \item 边频次：路径常见程度；
  \item 平均耗时：相邻活动之间的等待或处理时间；
  \item 低频高风险边：可能是例外路径或合规问题；
  \item 高频长耗时边：通常是流程改进候选。
\end{itemize}
\end{frame}

\begin{frame}{流程变体}
\begin{block}{Variant}
活动序列完全相同的 case 属于同一个流程变体。
\end{block}
\begin{itemize}
  \item 标准路径：制度流程的主要执行方式；
  \item 跳步路径：缺少某些活动；
  \item 返工路径：活动重复或出现 rework；
  \item 异常路径：活动顺序违反参考流程。
\end{itemize}
\small
变体分析回答：流程实际有多少种走法，哪些路径最常见、最慢、最风险。
\end{frame}

\begin{frame}{第10周事件日志}
\begin{center}
\begin{tabular}{lr}
\toprule
指标 & 数值\\
\midrule
事件数 & 67\\
case 数 & 8\\
流程变体数 & 4\\
平均吞吐时间 & 114.88 小时\\
返工 case & O2C003, O2C006\\
\bottomrule
\end{tabular}
\end{center}
\small
样例流程为 order-to-cash：接单、信用检查、拣货、质检、包装、发货、开票、收款。
\end{frame}

\begin{frame}{主要流程变体}
\begin{center}
\begin{tabular}{lrrr}
\toprule
变体 & case 数 & 占比 & 平均吞吐时间\\
\midrule
V01 标准路径 & 4 & 50.00\% & 102.75\\
V02 返工路径 & 2 & 25.00\% & 183.00\\
V03 反序异常 & 1 & 12.50\% & 73.00\\
V04 跳过信用检查 & 1 & 12.50\% & 69.00\\
\bottomrule
\end{tabular}
\end{center}
\small
返工路径的平均吞吐时间明显高于标准路径。
\end{frame}

\begin{frame}{Directly-Follows 结果}
\begin{center}
\begin{tabular}{lrr}
\toprule
边 & 频次 & 平均小时\\
\midrule
pack \(\rightarrow\) ship & 8 & 16.25\\
pick \(\rightarrow\) quality & 8 & 4.50\\
credit \(\rightarrow\) pick & 7 & 9.29\\
invoice \(\rightarrow\) payment & 7 & 71.57\\
order \(\rightarrow\) credit & 7 & 8.00\\
\bottomrule
\end{tabular}
\end{center}
\begin{block}{瓶颈候选}
开票到收款平均 71.57 小时，是最慢的相邻活动边。
\end{block}
\end{frame}

\begin{frame}{合规异常}
\begin{center}
\begin{tabular}{ll}
\toprule
case & 问题\\
\midrule
O2C004 & missing credit\_check\\
O2C007 & activity order violation\\
\bottomrule
\end{tabular}
\end{center}
\begin{itemize}
  \item 缺少信用检查可能是低风险快速通道，也可能是规则执行缺口；
  \item 发票早于发货可能是系统补录、特殊合同或真实流程反序；
  \item 异常解释必须回到业务规则和数据质量。
\end{itemize}
\end{frame}

\begin{frame}{数据质量检查}
\begin{itemize}
  \item case id 是否稳定，一张订单是否拆成多个履约单？
  \item timestamp 是开始时间、完成时间还是补录时间？
  \item activity 命名是否一致，粒度是否稳定？
  \item resource 是否表示实际执行者？
  \item 例外路径是违规、允许路径还是数据缺失？
  \item 流程发现是否结合订单金额和客户影响？
\end{itemize}
\end{frame}

\begin{frame}[fragile]{代码入口}
\begin{lstlisting}[language=bash]
PYTHONPATH=src python3 src/bdm_decision/cases/week10_process_mining.py
\end{lstlisting}

\begin{lstlisting}[language=Python]
from bdm_decision.process_mining.event_log import (
    load_event_log, summarize_cases, summarize_variants,
    directly_follows_graph
)

events = load_event_log()
cases = summarize_cases(events)
variants = summarize_variants(cases)
edges = directly_follows_graph(events)
\end{lstlisting}
\end{frame}

\begin{frame}{课堂讨论}
\begin{enumerate}
  \item 为什么平均吞吐时间不能单独说明流程健康？
  \item invoice 到 payment 的长等待应由哪个部门负责？
  \item O2C004 缺少信用检查一定是违规吗？
  \item 如何判断返工是质量控制有效还是流程浪费？
  \item 过程挖掘结果如何进入下一轮优化或自动化？
\end{enumerate}
\end{frame}

\begin{frame}{本讲小结}
\begin{itemize}
  \item 事件日志将流程转化为 case、activity 和 timestamp；
  \item trace 和 variant 揭示实际路径；
  \item directly-follows graph 展示相邻活动关系和等待；
  \item 过程挖掘包括发现、合规检查和流程增强；
  \item order-to-cash 案例暴露支付延迟、返工和异常路径；
  \item 流程结论必须结合时间戳质量、业务规则和客户影响。
\end{itemize}
\end{frame}

\end{document}
